\chapter{多因子与计量：从横截面信号到研究材料}
\label{ch:lower-multifactor}

\begin{itemize}
  \item 直接先修：上册第 6、10、13、14、16、17 章；诊断入口见第 \ref{route:multifactor} 节。
  \item 学习产出：交付一个能从干净环境复现、具有独立小面板 oracle、失败协议和限制报告的因子研究材料。
  \item 研究边界：本章验证研究方法，不把合成收益解释为真实 alpha、实盘业绩或生产经验。
\end{itemize}

\section{横截面目标与时间索引}

在决策时点 $t$，资产 $i$ 的信号 $x_{i,t}$ 必须属于合法信息集 $\mathcal I_t$；被解释变量是之后才实现的收益 $r_{i,t+1}$。最小横截面模型写成
\[
 r_{i,t+1}=a_t+b_t x_{i,t}+u_{i,t+1}.
\]
这里 $b_t$ 是当期信号与下一期收益的横截面线性关系，不自动是结构性风险价格，更不证明因果。若把 $x_{i,t+1}$ 回填到 $t$，即使回归代码完全正确，目标也已变成未来信息可见的伪问题。

对中心化信号，含截距 OLS 斜率为
\[
 \hat b_t=
 \frac{\sum_i(x_{i,t}-\bar x_t)(r_{i,t+1}-\bar r_{t+1})}
      {\sum_i(x_{i,t}-\bar x_t)^2}.
\]
本章两期手工面板的真实斜率分别为 $0.018$ 与 $0.022$。程序必须在读取收益前先验证 \texttt{signal\_date < return\_date}；相等或逆序都以稳定类别拒绝。

\subsection*{第一期逐行手算}

第一期信号为 $(-1,0,2,1)$，下一期收益为 $(-0.008,0.010,0.046,0.028)$。两列均值分别为 $0.5$ 与 $0.019$。因此
\[
 \sum_i(x_i-0.5)^2=5,
 \qquad
 \sum_i(x_i-0.5)(r_i-0.019)=0.09,
\]
斜率为 $0.09/5=0.018$，截距为 $0.019-0.018(0.5)=0.010$。第二期信号 $(-2,-1,2,1)$ 的平方离差和为 $10$，协方差分子为 $0.22$，故斜率为 $0.022$。这组算术是程序外的正确性来源。

\begin{longtable}{rrrr}
\toprule
时期 & 信号最小值 & 信号最大值 & 多空价差收益 \\
\midrule
1 & $-1$ & $2$ & $0.046-(-0.008)=0.054$ \\
2 & $-2$ & $2$ & $0.049-(-0.039)=0.088$ \\
\bottomrule
\end{longtable}

\begin{diagnostic}
若“修复”日期错误的方法是整体 \texttt{shift(-1)}，仍要逐列检查发布日期、财报修订、指数成分和停牌信息。数组对齐正确不等于时点数据正确。
\end{diagnostic}

\section{Fama--MacBeth 两阶段估计}

第一阶段对每个 $t$ 独立做横截面回归，得到 $\hat b_1,\ldots,\hat b_T$。第二阶段报告时间均值
\[
 \bar b=\frac1T\sum_{t=1}^T\hat b_t,
 \qquad
 \operatorname{se}(\bar b)
 =\sqrt{\frac{1}{T(T-1)}\sum_{t=1}^T(\hat b_t-\bar b)^2}.
\]
若 $\hat b_t$ 存在序列相关，应把第二式替换为声明带宽的 HAC 长期方差估计；若资产间误差相关，不能把同一期的公司数误当成独立时间样本。Fama--MacBeth 是估计与推断协议，不是“先回归两遍就稳健”\cite{fama1973risk}。

固定小面板给出 $\bar b=0.020$、时间序列标准误 $0.002$。这两个数来自两期斜率的手算，而不是从被测程序重新生成期望答案。

具体地，两个斜率相对均值的偏差为 $(-0.002,0.002)$，样本方差为 $8\times10^{-6}$；均值方差再除以 $T=2$，得到 $4\times10^{-6}$，开方即 $0.002$。真实研究的 $T$ 应远大于 2，本例只用于让读者逐项核对公式。月份重叠、滚动收益或平滑信号会令 $\hat b_t$ 自相关，此时必须保存完整的系数时间序列，才能复核 HAC 带宽与有限样本修正。

\section{暴露、中性化与病态矩阵}

令 $Z_t$ 包含截距、规模和行业暴露。信号中性化是投影残差
\[
 x_t^{\perp}=x_t-Z_t(Z_t^TZ_t)^{-1}Z_t^Tx_t,
 \qquad Z_t^Tx_t^{\perp}=0.
\]
实现不应显式求逆，而应使用 QR、SVD 或最小二乘求解，并在条件数超过声明阈值时拒绝结果。本章 oracle 检查截距、规模和行业约束的最大绝对残差；另用完全共线的规模与行业列触发病态矩阵失败。

第一期采用规模 $(-1,-1,1,1)$ 和行业虚拟变量 $(0,1,0,1)$。最小二乘拟合系数为 $(0.5,1,0)$，中性化残差是
\[
 x^\perp=(-0.5,0.5,0.5,-0.5).
\]
残差和为零，与规模及行业列的内积也都为零；同时残差并非零向量，因此该例既验证约束，又保留一个可以继续排序的信号。若设计矩阵把同一行业同时编码成全套虚拟变量和截距，就会出现 dummy-variable trap；此时求得某个数值解不等于模型已识别。

中性化改变了信号的统计目标。若规模既是风险暴露又承载真实预测信息，残差化会删除两者；“更中性”不是无条件更优。生产研究还需事先确定行业分类版本、缺失暴露处理和极端值规则。

\section{IC、Rank IC、衰减与多重检验}

Pearson IC 是 $x_{i,t}$ 与 $r_{i,t+1}$ 的横截面相关；Rank IC 是两者秩的相关。前者对线性尺度和极端值敏感，后者只保留次序且需声明并列秩规则。本章精确线性小面板二者都等于 $1$，这只是 oracle 设计，不代表真实市场会出现完美相关。

衰减分析必须事先确定预测跨度 $h$：
\[
 \operatorname{IC}(h)=\operatorname{Corr}_i(x_{i,t},r_{i,t\to t+h}).
\]
每尝试一个信号、窗口、行业处理、收益跨度或异常值规则，就消耗一次研究选择权。固定随机种子后，程序生成 60 期独立标准正态收益和 20 个独立标准正态信号；这是真正的零信号搜索，而不是预先写好的 $p$ 值。Fisher $z$ 近似产生的 20 个 $p$ 值中，朴素 $p<0.05$ 会留下 2 项；Benjamini--Hochberg 阶梯在 $5\%$ 水平拒绝 0 项。报告必须同时保存随机种子、总尝试数、筛选规则与未入选结果\cite{white2000reality,bailey2017backtest}。

对事先确定零信号搜索排序后，最小三个 $p$ 值约为 $0.027691,0.037847,0.073846$。BH 前三个阈值分别为 $0.0025,0.005,0.0075$，均未通过；后续值也不满足阶梯条件，所以拒绝数为 0。这里的“0 项”不是证明所有信号都无效，而是说明在当前检验族和错误率目标下证据不足。若研究者删掉 18 个失败尝试再把 $m$ 写成 2，检查应直接判定报告不可审计。

第二个事先确定面板覆盖 2016--2024 年，得到
\[
\begin{aligned}
(\mathrm{IC}_{\mathrm{train}},\mathrm{IC}_{\mathrm{select}},\mathrm{IC}_{\mathrm{test}})
  &=(0.583807,0.711369,0.726531),\\
(\mathrm{IC}_{1},\mathrm{IC}_{3},\mathrm{IC}_{6})
  &=(0.646657,0.445136,0.186461).
\end{aligned}
\]
这些数值来自固定种子的合成数据，作用是验证时间切分和衰减报告链路，不是市场经验结论。

\section{成本、容量与可交易性}

把最高信号资产做多、最低信号资产做空，两个时期的手工价差收益分别为 $0.054$ 与 $0.088$，故毛收益均值为 $0.071$。若一次往返摩擦为 $0.004$，则
\[
 R^{\mathrm{net}}=0.071-0.004=0.067.
\]
当资金规模使额外冲击增加 $0.006$ 时，净收益降为 $0.061$。这只是说明容量必须进入结果函数；真实冲击应由参与率、成交量、波动率和执行方式估计，不能用一个常数永久外推。

一个合格结果表至少并列报告：毛收益、佣金、价差、冲击、借券、净收益、双边换手、资金规模和可成交比例。若只报告 Sharpe 或 IC，读者无法从信号质量恢复可交易性。容量也不是单个“最大 AUM”：应画出资金规模到净绩效、参与率和成交完成率的曲线，并对冲击模型参数做敏感性分析。

\begin{implementationnote}
执行下面的命令：
\begin{lstlisting}[language=bash]
uv run python notebooks/lower/ch01_multifactor.py \
  evidence/lower-ch01/oracle.json
\end{lstlisting}
输出同时包含系数、IC、中性化约束、多重检验、毛净收益与两种失败检查；任一字段漂移都会非零退出。
\end{implementationnote}

\subsection*{A 股制度边界}
涨跌停、停牌、T+1、融券可得性和复权数据发布时间会改变可交易集合或买卖方向，因此必须进入时点资产池与交易台账。除此之外，横截面估计、FDR 和条件数规则仍按通用数学表述，不把市场特例写成普遍定理。

\section{本路线的综合项目：全球因子研究材料}

Capstone 需事先确定研究问题、时点数据协议、基准信号、样本外窗口、尝试次数、成本容量模型和限制报告；必须包含一个因未来暴露或病态中性化矩阵而被研究设计拒绝的失败案例。交付只能声称“该研究材料在声明数据与协议下可复现”，不能声称稳定盈利、生产部署或团队经验。

配套材料把数据、源码、独立小面板和报告分开保存；读者需要看到的是每一部分承担的职责，而不是目录名称。另一位研究者应能从同一输入重算系数表、成本分解与限制报告。

本章报告说明训练、选择和最终测试窗口互不重叠，衰减滞后固定为 1、3、6 个月，并把全球核心设定与 A 股制度边界分开陈述。它是本章可复现研究材料的一部分，而不是留给未来综合项目的空白模板。

\section{四级问题}
\begin{enumerate}[leftmargin=*]
  \item \textbf{口述概念：}横截面预测系数、风险价格和因果效应为何不是同一概念？
  \item \textbf{口述概念：}IC 与 Rank IC 各保留什么信息，何时会给出不同排序？
  \item \textbf{笔试推导：}从两期斜率 $0.018,0.022$ 推导 Fama--MacBeth 均值和标准误。
  \item \textbf{笔试推导：}证明投影残差满足 $Z^Tx^\perp=0$，并说明秩亏时哪里失效。
  \item \textbf{数值编程：}扩展 fixture 至三期，同时保持独立斜率和价差收益台账。
  \item \textbf{数值编程：}篡改 signal date、行业列和一个 oracle 字段，记录三个稳定失败类别。
  \item \textbf{研究判断：}500 个因子中 20 个通过 BH，怎样证明其中任何一个具有成本后样本外价值？
  \item \textbf{研究判断：}一个全球因子在 A 股失效，区分数学失效、数据失效和制度失效。
\end{enumerate}

